\begin{ouverture}[probabilites]{20}{Comprendre le hasard}{Partie 4 · Données, hasard et instructions}
{Chances · Équiprobabilité · Fraction · Fréquence}
\objectif{Savoir qu’une probabilité est un nombre compris entre 0 et 1.}
\objectif{Calculer une probabilité dans une situation d’équiprobabilité.}
\objectif{Comparer les résultats d’essais répétés à une probabilité calculée.}
\prerequis{« a chances sur b » au CM2 ; fractions, décimaux, pourcentages.}
\end{ouverture}
\pagestyle{otMain}\essentiel\fontsize{11}{15}\selectfont
\phantomsection\label{cours-20-1}
\essSection{1}{Du langage des chances au nombre}
\begin{definition}Une expérience est aléatoire quand on connaît les résultats possibles, appelés issues, sans pouvoir prévoir lequel se produira. La probabilité d’un événement est un nombre compris entre 0 et 1 : 0 pour un événement impossible, 1 pour un événement certain.\end{definition}
\label{fig-cours-20-1-cours-echelle-proba}\figureGeo{cours-echelle-proba}

\essSection{2}{Recenser les possibilités et calculer}
\begin{propriete}Quand toutes les issues ont la même chance de se produire, on parle d’équiprobabilité. La probabilité d’un événement est alors le quotient du nombre d’issues favorables par le nombre total d’issues : « $a$ chances sur $b$ » devient le nombre $\frac ab$.\end{propriete}
\begin{exemple}Avec un dé équilibré à six faces, la probabilité d’obtenir 6 est $\frac16$ ; celle d’obtenir un nombre pair est $\frac36=\frac12$, car 2, 4 et 6 conviennent.\end{exemple}
\clearpage
\phantomsection\label{cours-20-2}
\essSection{3}{Relier les écritures et comparer}
\begin{savoirfaire}[Changer d’écriture]Une probabilité peut s’écrire avec une fraction, un décimal ou un pourcentage : $\frac12=0{,}5=50$~\% ; $\frac14=0{,}25=25$~\%. Pour comparer deux probabilités, on les écrit sous la même forme.\end{savoirfaire}
\begin{savoirfaire}[Recenser deux étapes]Pour deux pièces lancées, un arbre ou un tableau recense les 4 issues PP, PF, FP et FF, qui sont équiprobables : obtenir une pile et une face a donc 2 chances sur 4.\end{savoirfaire}
\essSection{4}{Observer des essais répétés}
\begin{definition}Quand on répète une expérience, la fréquence d’un résultat est le quotient de son nombre d’apparitions par le nombre d’essais.\end{definition}
\label{fig-cours-20-2-cours-frequences}\figureGeo{cours-frequences}

\begin{propriete}Quand le nombre d’essais devient grand, la fréquence se rapproche en général de la probabilité. On ne peut pourtant jamais prévoir le résultat d’un essai : le hasard n’a pas de mémoire.\end{propriete}

\clearpage\rubrique{Découverte}
\begin{blocLivre}{ex-20-1}\begin{exercice}[1]{De l’impossible au certain}
On lance un dé équilibré à six faces. Donner la probabilité de chaque événement.\begin{questions}
\item Obtenir 7.
\item Obtenir un nombre inférieur ou égal à 6.
\item Obtenir 6.
\item Obtenir un nombre pair.
\end{questions}
\end{exercice}\end{blocLivre}
\begin{blocLivre}{ex-20-2}\begin{exercice}[1]{Une urne}
Une urne contient des boules indiscernables au toucher : 3 rouges, 5 bleues et 2 vertes. On tire une boule au hasard. Calculer la probabilité de tirer une boule rouge, bleue, verte, puis une boule qui n’est pas rouge.
\label{fig-ex-20-2}\figureGeo{hasard-2}
\end{exercice}\end{blocLivre}
\begin{blocLivre}{ex-20-3}\begin{exercice}[1]{Changer d’écriture}
Écrire chaque probabilité sous les deux autres formes : $\frac14$ ; $\frac35$ ; $0{,}7$ ; 15~\%.
\end{exercice}\end{blocLivre}
\begin{blocLivre}{ex-20-4}\begin{exercice}[1]{La roue}
La roue est partagée en 8 secteurs identiques numérotés de 1 à 8. On la fait tourner. Calculer la probabilité d’obtenir un nombre pair, un nombre supérieur ou égal à 6, puis un multiple de 3.
\label{fig-ex-20-4}\figureGeo{hasard-4}
\end{exercice}\end{blocLivre}
\begin{blocLivre}{ex-20-5}\begin{exercice}[1]{Une pièce sans mémoire}
Une pièce équilibrée vient de tomber trois fois de suite sur pile. Lucas pense que face a plus de chances de sortir au quatrième lancer. A-t-il raison ?
\end{exercice}\end{blocLivre}

\rubrique{Entraînement}
\begin{blocLivre}{ex-20-6}\begin{exercice}[2]{Deux pièces}
On lance deux pièces équilibrées. L’arbre recense les issues. Quelle est la probabilité d’obtenir deux piles ? Une pile et une face ?
\label{fig-ex-20-6}\figureGeo{hasard-6}
\end{exercice}\end{blocLivre}
\begin{blocLivre}{ex-20-7}\begin{exercice}[2]{Deux dés}
On lance deux dés équilibrés et on additionne les résultats. Le tableau recense les 36 issues. Quelle est la probabilité d’obtenir une somme égale à 7 ? Égale à 12 ?
\label{fig-ex-20-7}\figureGeo{hasard-7}
\end{exercice}\end{blocLivre}
\begin{blocLivre}{ex-20-8}\begin{exercice}[2]{Fréquence et probabilité}
Emma lance un dé 200 fois et obtient 36 fois le six. Calculer la fréquence du six, puis la comparer à la probabilité.
\end{exercice}\end{blocLivre}
\begin{blocLivre}{ex-20-9}\begin{exercice}[2]{Deux urnes}
L’urne A contient 2 boules rouges sur 5 boules ; l’urne B, 3 boules rouges sur 8. Dans quelle urne a-t-on le plus de chances de tirer une boule rouge ?
\label{fig-ex-20-9}\figureGeo{hasard-9}
\end{exercice}\end{blocLivre}
\begin{blocLivre}{ex-20-10}\begin{exercice}[2]{Une roue inégale}
La roue comporte trois secteurs A, B et C. Tom dit : « Il y a trois secteurs, donc chacun a une chance sur trois. » A-t-il raison ?
\label{fig-ex-20-10}\figureGeo{hasard-10}
\end{exercice}\end{blocLivre}

\rubrique{Problèmes}
\begin{blocLivre}{ex-20-11}\begin{exercice}[3]{Le jeu de 32 cartes}
On tire une carte au hasard dans un jeu de 32 cartes : 4 couleurs (cœur, carreau, trèfle, pique) de 8 cartes, dont un roi par couleur. Quelle est la probabilité de tirer un cœur ? Un roi ? Le roi de cœur ?
\end{exercice}\end{blocLivre}
\begin{blocLivre}{ex-20-12}\begin{exercice}[3]{Le sac de bonbons}
Un sac contient 20 bonbons. La probabilité de tirer un bonbon à la fraise est $0{,}35$. Combien y a-t-il de bonbons à la fraise ? On ajoute 5 bonbons à la fraise : quelle est la nouvelle probabilité, en pourcentage ?
\end{exercice}\end{blocLivre}
\begin{blocLivre}{ex-20-13}\begin{exercice}[3]{La punaise}
Une punaise lancée retombe pointe en haut ou pointe en bas ; les deux issues n’ont pas forcément la même chance. Une classe a réalisé $1\,000$ lancers et obtenu 620 fois « pointe en haut ». La courbe montre l’évolution de la fréquence. Quelle probabilité peut-on proposer ?
\label{fig-ex-20-13}\figureGeo{hasard-13}
\end{exercice}\end{blocLivre}
\begin{blocLivre}{ex-20-14}\begin{exercice}[3]{Un jeu équitable ?}
On lance deux dés et on additionne. Alice gagne si la somme vaut 7 ; Bob gagne si la somme vaut 2, 3 ou 12. Le jeu est-il équitable ? Ajouter à Bob une somme pour le rendre équitable.
\end{exercice}\end{blocLivre}
\begin{blocLivre}{ex-20-15}\begin{exercice}[3]{Vrai ou faux}
Nina affirme : « Avec deux dés, obtenir 6 avec le premier dé et 6 avec le second est moins probable qu’obtenir 3 avec le premier et 5 avec le second. » A-t-elle raison ?
\end{exercice}\end{blocLivre}

\clearpage
